\begin{figure}[t]
\centering
\begin{tikzpicture}[x=1.35cm,y=1cm,>=Stealth,
    every node/.style={font=\small}]
  \fill[blue!4] (-4,-1.4) rectangle (-2,1.4);
  \fill[blue!4] (2,-1.4) rectangle (4,1.4);
  \draw[dashed,gray] (-2,-1.55)--(-2,1.6);
  \draw[dashed,gray] (2,-1.55)--(2,1.6);
  \draw[gray] (0,-1.55)--(0,1.6);
  \draw[line width=2pt,red!55!black] (0,-.9)--(0,.9);
  \draw[thick] (-2,.6)--(2,.6)
    .. controls (4.25,.6) and (4.25,-.6) .. (2,-.6)
    --(-2,-.6)
    .. controls (-4.25,-.6) and (-4.25,.6) .. (-2,.6);
  \fill[white] (-3.6875,0) circle (.10);
  \draw[thick] (-3.7875,0)--(-3.5875,0);
  \fill[white] (3.6875,0) circle (.10);
  \draw[thick] (3.5875,0)--(3.7875,0);
  \fill (0,.6) circle (.055);
  \fill (0,-.6) circle (.055);
  \node[above left] at (-2,.6) {$+$};
  \node[below left] at (-2,-.6) {$+$};
  \node[above right] at (2,.6) {$-$};
  \node[below right] at (2,-.6) {$-$};
  \node[above right] at (0,.6) {$j$};
  \node[below right] at (0,-.6) {$i$};
  \node[right,red!55!black] at (0,0) {$I$};
  \node at (-3,-1.15) {left cap};
  \node at (3,-1.15) {right cap};
  \node[above] at (-2,1.6) {$a$};
  \node[above] at (0,1.6) {$c$};
  \node[above] at (2,1.6) {$b$};
  \node[above] at (-3.55,.86) {source};
  \node[above] at (3.55,.86) {source};
  \draw[<->] (-1.85,-1.15)--(-.15,-1.15);
  \draw[<->] (.15,-1.15)--(1.85,-1.15);
  \node[below] at (-1,-1.15) {$t_-$};
  \node[below] at (1,-1.15) {$t_+$};
\end{tikzpicture}
\caption{A same-loop contribution to the insertion. The selected cap arches
carry two sources: the endpoint signs are \(++\) at the left seam and
\(--\) at the right seam. The two sinks \(i,j\) lie in the middle
interval \(I\). The insertion changes the permitted color assignments
while preserving the uncolored loop. The positive propagation widths
\(t_-\) and \(t_+\) control the sum over mode states.}
\label{fig:two-particle-insertion}
\end{figure}
